\section{One Clock, Two Proofs}

A burden enters when a reading asks to be resolved. The word burden is useful
because it does not pretend that resolution is free. A finite grammar must
carry the question, test the admissible routes by which the question could
close, and record the cost of closure. If there is no permitted route, the
burden remains unresolved in that grammar. If there is a permitted route,
the grammar must still say how much clock was spent to reach it.

The single clock from the foundation returns here in a sharper role. A clock
is not an external river in which arguments drift. It is the grammar's
bounded count of permitted turns. It says when a trial begins, when an
admissible update has occurred, and when the record has spent enough
distinguishable steps to accept a result. A result without a clock would
have no elapsed shape. A clock without a result would count only motion. A
measured burden needs both.

The first reading of the burden is decidability. Decidability asks whether
the grammar can resolve the burden at all. It is a yes-or-no reading, but it
is not the primitive truth collapse of Chapter 01 repeated at a higher
register. The grammar now has a needle, a finite gauge, and a positive
invariant. It can ask a disciplined question: under these permitted rules,
with these boundaries and this finite record, is there a route by which the
burden closes? If the route closes, the burden is decidable for that
grammar. If it cannot close, the grammar is not licensed to treat nonclosure
as false; it must read nonclosure as unresolved or obstructed relative to the
available sensor.

The second reading is elapsed cost. Elapsed cost asks how many permitted
turns the resolution spends. Two burdens can both resolve and still differ
in cost. One may close directly because the boundary already carries the
needed identification. Another may require refinements, tests, and
quotients before the same kind of record can be accepted. The answerability
of a burden and the time spent answering it are therefore two readings of one
burden, not two unrelated facts.

This distinction matters because the grammar is finite. If decidability were
read without cost, every resolved burden would look equally cheap. If cost
were read without decidability, the grammar would accumulate steps without
knowing whether the steps had closed anything. Measurement needs the pair.
It must know that a burden can be resolved, and it must know how the clock
was spent in resolving it.

The phrase ``one clock, two proofs'' should be heard in this sense. A proof
is not a private certificate floating above the record. It is a route of
resolution that the grammar can read. One route may emphasize the fact of
closure: the burden admits a terminating reading. Another may emphasize the
elapsed work: the same burden consumes a calibrated amount of clock. The two
routes identify the same burden only when the grammar can compare them under
one clock. Otherwise they might be two incomparable stories about different
records.

Tail latency supplies a useful illustration. A networked measurement can
have a local answer ready while a boundary reconciliation still waits for
the slowest connected interface. The grammar of the record does not accept
the answer merely because one path was quick. It accepts the answer when the
connected boundary has been reconciled. The elapsed cost is therefore read
at the slowest admissible closure, while decidability reads whether closure
is available at all. The example is not the argument; it merely shows how a
single clock can expose two different aspects of the same burden.

The same pairing appears in a calibration chain. A sensor may be able to
decide whether a reading lies inside tolerance. That is the decidability
face. But the number of repetitions required before the tolerance is
accepted is the elapsed-cost face. If the repetitions can be compressed by a
better boundary commitment, the burden has the same answer but a different
cost. If no admissible repetition can settle the tolerance, the cost is not
large; the burden is unresolved for that grammar.

The operator enters because the burden needs a concrete reader. A bare
declaration that two routes agree would only rename the problem. The grammar
must carry a finite object on which closure and cost can both be read. That
object must have places where differences are tested, a boundary where drift
is forbidden, and an energy reading that knows when activity has vanished.
Without such a reader, the clock can count turns but cannot say what the
turns have resolved.

Thus the first beat of the chapter grounds the invariant in a burden read
twice. Decidability reads answerability. Elapsed cost reads the clock spent
by answerability. The single clock identifies the two readings as readings
of one burden rather than rival accounts. Once this is required, the next
question is forced. What finite reader can make zero burden visible as zero
activity rather than as an invisible drift? The answer is the anchored
chain.
